% ---------------------------------------------------------------------------
\section{Introduction}
\label{sec:intro}

Reinforcement learning with verifiable rewards (RLVR) is now a central ingredient in training reasoning models \citep{tulu3,deepseekmath,deepseekr1,qwen3}.
The recipe is simple. Sample several solutions to a problem, reward those whose final answer matches the reference, and raise the probability of rewarded solutions with a policy gradient such as group relative policy optimization (GRPO; \citealp{deepseekmath}).
Its bottleneck is equally simple: every training problem needs a trusted reference answer.
Such answers are expensive to produce, and they become harder to obtain as models approach expert ability.

This has motivated \emph{label-free} or \emph{unsupervised} RLVR, in which the model rewards itself.
The most common choice is majority voting \citep{ttrl,shafayat2025srt}. The answer that appears most often among the $K$ samples is treated as a pseudo-label, and samples that agree with it are rewarded.
Related intrinsic signals include self-certainty and entropy \citep{intuitor,agarwal2025entmin,rent,empo}.
These methods often help early in training, but with prolonged training they \emph{collapse}.
The self-assessed reward keeps rising while true accuracy falls, and in the extreme the model emits the same answer for every problem \citep{shafayat2025srt,he2026urlvr}.

A striking feature of this collapse is its dependence on the \emph{number of training problems}.
In a systematic study, \citet{he2026urlvr} found that training on $\le128$ problems was stable, whereas ``larger datasets ($\ge512$) consistently exhibit reward hacking''.
The same algorithm is safe with few problems and fails with many.
This is the signature of a \emph{collective} phenomenon, and it is not explained by existing theory.
The theory of intrinsic rewards is single-problem: each problem is sharpened toward its initial majority answer \citep{huang2025sharpening,he2026urlvr}.
Such sharpening helps when the initial majority is correct and hurts when it is wrong, but it says nothing about why \emph{all} problems should be dragged toward a common answer, or why that should depend on how many there are.

\paragraph{This paper.}
We argue that dataset-size-dependent collapse is an instability of a \emph{shared direction} in parameter space.
Problems are coupled because they share network parameters \citep{jacot2018ntk,ren2025learning,fan2026sharpening}.
Any parameter that moves all answers the same way (a format, a preferred value, an output bias) collects the majority-vote signal from every problem at once, whereas each problem-specific parameter collects only its own share.
Our contributions are:
\begin{itemize}\setlength\itemsep{1pt}
\item \textbf{Mechanism.} In a solvable many-problem model, a shared perturbation of the state in which different problems prefer different answers grows, at each instant along the sharpening trajectory, exactly when the average slope $\langle F'\rangle$ of the single-problem self-reinforcement drift $F$ is positive. For $K=3$ samples per problem this holds when problems are still uncertain, with answer logits $|z|<\ln(2+\sqrt3)\approx1.32$, a Curie--Weiss-type paramagnetic instability (\cref{sec:theory}).
\item \textbf{A critical dataset size.} Local sharpening, slowed by $1/N$ for $N$ training problems, races against the shared mode, which grows at a rate $\propto\lambda^2$ independent of $N$, where $\lambda$ is the strength of the shared coupling. The result is a closed-form critical size $N_c\simeq\ln(1/\delta)/(\lambda^2\ln[F^\star/F(a)])$, where $\delta$ is the noise that seeds the shared mode, $a$ is the initial confidence and $F^\star$ is the drift at the edge of the uncertain range. It makes four falsifiable predictions.
\item \textbf{Verification.} Stochastic simulations with actual sampling, majority voting and minibatch stochastic gradient descent (SGD) confirm all four predictions. The transition is sharp; $N_c\propto\lambda^{-2}$ up to a logarithmic correction, with a single critical value $\Lambda_c\approx2.6$ of the amplification $\lambda^2N\ln[F^\star/F(a)]$ that also captures the dependence on initial confidence; batch size enters only logarithmically; and confident initial models never collapse, even under much longer training. The same law holds for GRPO, larger groups, many answers, heterogeneous confidence and Adam, the last with a larger constant (\cref{sec:sims,sec:robust}).
\item \textbf{A kernel version for networks.} Linearizing a network gives a version of the theory in terms of the empirical neural tangent kernel (NTK). It reveals a second, protective coupling, the generalization of sharpening across problems, under which adding data can stop causing collapse (\cref{sec:kernel}).
\item \textbf{A small Transformer.} In label-free RL, removing the shared direction by an antisymmetric readout prevents collapse entirely, adding it back causes collapse, and collapse grows with training-set size and base-model entropy. Predictions from the initial kernel rank runs by their risk of collapse, but they are poorly calibrated: they underestimate collapse and agree with the outcome no more often than always predicting collapse. Measuring the kernel during training identifies one source of this error, the erosion of the protective coupling by feature learning (\cref{sec:nn}).
\item \textbf{Stabilizing the uniform mode, and its limit.} A batch-level marginal-entropy bonus with weight $\beta$ damps the uniform shared mode. Its critical weight $\beta_c\le\tfrac12$ is derived from the theory and confirmed in simulation and in the Transformer. In the network, however, the instability then moves to the next collective kernel mode, so removing the imbalance of the answer marginal does not prevent a task-unrelated collective collapse, and accuracy is not preserved.
\end{itemize}
After reviewing related work, we analyze a single problem in isolation (\cref{sec:single}) and then couple many problems through a shared parameter (\cref{sec:theory}).
We test the resulting predictions in stochastic simulations (\cref{sec:sims}, with further robustness checks in \cref{sec:robust}), carry them over to a small Transformer (\cref{sec:nn}), and close with implications and limitations (\cref{sec:discussion}).

% ---------------------------------------------------------------------------
\section{Related Work}
\label{sec:related}

\paragraph{Label-free RL for reasoning.}
Test-time reinforcement learning (TTRL; \citealp{ttrl}) uses majority-vote pseudo-rewards on unlabeled test problems.
Self-training with majority votes \citep{shafayat2025srt}, self-certainty \citep{intuitor}, entropy minimization \citep{agarwal2025entmin,rent} and semantic-entropy objectives \citep{empo} are close relatives. Even random or spurious rewards can help some base models \citep{spurious}.
Several works report eventual reward hacking or collapse and propose remedies such as cross-view or cross-model rewarding \citep{corewarding}.
\citet{he2026urlvr} study unsupervised RLVR systematically. They show that intrinsic methods converge toward sharpening the model's initial distribution, that collapse depends on dataset size, and that base models with higher initial entropy fare worse.
We supply a mechanism for the dataset-size dependence and a quantitative criterion.

\paragraph{Sharpening and entropy collapse.}
\citet{huang2025sharpening} formalize self-improvement as sharpening toward high-likelihood responses.
\citet{cui2025entropy} analyze entropy collapse in RL for reasoning models through the covariance between log-probabilities and advantages.
These analyses treat a single prompt, or an aggregate over prompts, without the cross-problem coupling that drives our transition.

\paragraph{Coupling through shared parameters.}
The NTK \citep{jacot2018ntk} and the learning-dynamics view of fine-tuning \citep{ren2025learning} describe how an update on one example moves predictions on another.
Closest to us, \citet{fan2026sharpening} show for \emph{labeled} RLVR that finite-batch updates bias learning toward sampled modes and that collapse propagates through semantic coupling. Their analysis is a one-step linearization for a pair of queries.
We study the collective, many-problem dynamics of \emph{self}-rewarded RL and derive a size threshold.

\paragraph{Statistical physics of self-training and recursive training.}
\citet{takahashi2026selftraining} analyzes iterative self-training of linear classifiers with the replica method.
Model collapse under recursive training on generated data \citep{shumailov2024collapse} is a related but distinct failure that operates across model generations.
Our analysis uses mean-field and linear-stability tools \citep{nishimori2001}. It concerns on-policy sampling with majority-vote policy gradients, where the reward depends on the policy itself.
